\documentclass[12pt, a4paper]{article}

% ── PACKAGES ──────────────────────────────────────────────────
\usepackage[utf8]{inputenc}
\usepackage[T1]{fontenc}
\usepackage{geometry}
\geometry{margin=2.5cm}
\usepackage{hyperref}
\hypersetup{
    colorlinks=true,
    linkcolor=blue,
    urlcolor=blue,
    citecolor=blue,
    pdftitle={ILC as the Geometric Inverse of TNBC:
              Structural Lock vs Epigenetic Lock on the
              Same Luminal Identity Axis},
    pdfauthor={Eric Robert Lawson},
    pdfkeywords={ILC, invasive lobular carcinoma,
                 TNBC, triple-negative breast cancer,
                 FOXA1, EZH2, CDH1, geometric inverse,
                 structural lock, epigenetic lock,
                 luminal identity axis, fulvestrant,
                 tazemetostat, Waddington landscape,
                 attractor geometry, OrganismCore,
                 GSE176078, breast cancer, Type 3,
                 Type 2, E-cadherin, PRC2, H3K27me3}
}
\usepackage{amsmath}
\usepackage{booktabs}
\usepackage{longtable}
\usepackage{array}
\usepackage{parskip}
\usepackage{microtype}
\usepackage{authblk}
\usepackage{titlesec}
\usepackage{fancyhdr}
\usepackage{xcolor}
\usepackage{float}
\usepackage{enumitem}
\usepackage{setspace}
\usepackage{multirow}

% ── COLOURS ───────────────────────────────────────────────────
\definecolor{repobox}{RGB}{240,247,255}
\definecolor{findingbox}{RGB}{245,255,245}
\definecolor{convergentbox}{RGB}{255,250,235}
\definecolor{mechanismbox}{RGB}{250,245,255}
\definecolor{novelbox}{RGB}{255,248,220}
\definecolor{ilcbox}{RGB}{240,255,240}
\definecolor{tnbcbox}{RGB}{255,240,240}
\definecolor{axisbox}{RGB}{245,245,255}

% ── HEADER / FOOTER ───────────────────────────────────────────
\pagestyle{fancy}
\fancyhf{}
\lhead{\footnotesize CS-LIT-7: ILC as the Geometric
       Inverse of TNBC}
\rhead{\footnotesize OrganismCore \textbar{} 2026-03-06}
\rfoot{\small Page \thepage}
\renewcommand{\headrulewidth}{0.4pt}
\setlength{\headheight}{24pt}
\addtolength{\topmargin}{-10pt}

% ── SECTION FORMATTING ────────────────────────────────────────
\titleformat{\section}
  {\large\bfseries}{\thesection}{1em}{}[\titlerule]
\titleformat{\subsection}
  {\normalsize\bfseries}{\thesubsection}{1em}{}
\titleformat{\subsubsection}
  {\small\bfseries}{\thesubsubsection}{1em}{}

% ══════════════════════════════════════════════════════════════
% TITLE
% ══════════════════════════════════════════════════════════════
\title{
    \vspace{-1em}
    \textbf{ILC as the Geometric Inverse of TNBC:}\\[0.5em]
    \large Structural Lock versus Epigenetic Lock on
    the Same Luminal Identity Axis --- Mechanistically
    Opposite Entry Points for Therapy\\[0.4em]
    \normalsize \textit{OrganismCore Framework ---
    Document CS-LIT-7}\\[0.3em]
    \small Derived from Waddington Attractor Geometry
    of 19,542 Single Cancer Cells (GSE176078)\\[0.3em]
    \small Confirmed by ILC Biology (Frontiers Oncology
    2025) and TNBC Epigenetics (Schade et al.\
    \textit{Nature} 2024)
}

\author[1]{Eric Robert Lawson}
\affil[1]{
    OrganismCore\\
    \href{mailto:OrganismCore@proton.me}
         {OrganismCore@proton.me}\\
    ORCID:
    \href{https://orcid.org/0009-0002-0414-6544}
         {0009-0002-0414-6544}
}

\date{March 6, 2026}

% ══════════════════════════════════════════════════════════════
\begin{document}
% ══════════════════════════════════════════════════════════════

\maketitle
\thispagestyle{fancy}

% ── REPOSITORY POINTER ────────────────────────────────────────
\begin{center}
\colorbox{repobox}{
\begin{minipage}{0.88\textwidth}
\vspace{0.5em}
\centering
\textbf{Full computational record --- scripts, data
caches, reasoning artifacts, and raw logs:}\\[0.4em]
\href{https://github.com/Eric-Robert-Lawson/attractor-oncology}
{\texttt{https://github.com/Eric-Robert-Lawson/attractor-oncology}}
\\[0.3em]
\small Repository ID: 1172048282 \textbar{}
MIT License \textbar{}
Python \textbar{}
Public\\[0.2em]
Part of the OrganismCore BRCA Cross-Subtype Analysis
(BRCA-S8h, 2026-03-05). All predictions locked from
attractor geometry \textbf{before} any literature review
was performed.
\vspace{0.5em}
\end{minipage}}
\end{center}

\vspace{0.8em}

% ── CENTRAL FINDING BOX ───────────────────────────────────────
\begin{center}
\colorbox{axisbox}{
\begin{minipage}{0.88\textwidth}
\vspace{0.5em}
\centering
\textbf{THE GEOMETRIC INVERSION --- STATED ONCE:}\\[0.6em]
\begin{tabular}{p{3.5cm}p{4cm}p{4cm}}
\toprule
& \textbf{\textcolor{green!50!black}{ILC}} &
  \textbf{\textcolor{red!70!black}{TNBC}} \\
\midrule
FOXA1 &
Hyperactivated \textit{above} normal LumA &
Epigenetically silenced \textit{below} any
normal reference \\[0.3em]
EZH2 &
Low &
$+189\%$ above normal (highest of all subtypes) \\[0.3em]
CDH1 (E-cadherin) &
Absent (${\sim}65\%$ mutation,
${\sim}35\%$ methylation) &
Present \\[0.3em]
Lock type &
Structural (CDH1 loss) &
Epigenetic (EZH2/PRC2 H3K27me3) \\[0.3em]
Attractor geometry &
Type~3: Structural lock &
Composite Type~1$\rightarrow$2:
Epigenetic lock \\[0.3em]
Therapeutic entry &
Full ER degradation
(fulvestrant $>$ AI) &
EZH2 inhibition first
(tazemetostat), then fulvestrant \\
\bottomrule
\end{tabular}
\vspace{0.5em}
\end{minipage}}
\end{center}

\vspace{0.8em}

% ══════════════════════════════════════════════════════════════
\begin{abstract}
% ══════════════════════════════════════════════════════════════
\noindent
Invasive Lobular Carcinoma (ILC) and Triple-Negative
Breast Cancer (TNBC) are typically described as
opposite extremes of the breast cancer spectrum ---
ILC is ER-positive, low-grade, and indolent; TNBC
is ER-negative, high-grade, and aggressive. The
OrganismCore Waddington attractor geometry framework,
applied to 19,542 single cancer cells (GSE176078),
reveals that this clinical opposition is the surface
expression of a precise geometric inversion: both
subtypes are arrested forms of luminal identity, but
by mechanistically opposite locks acting on the same
axis.

ILC is a \textbf{structural lock} (Attractor Type~3):
FOXA1 is hyperactivated above normal luminal levels,
driving an amplified ER circuit; CDH1 (E-cadherin)
is absent (${\sim}65\%$ mutation, ${\sim}35\%$
methylation), removing the adhesion constraint that
normally organises luminal cells; EZH2 is low.
The cell is over-committed to luminal identity
with no structural containment.

TNBC is an \textbf{epigenetic lock}
(Composite Type~1$\rightarrow$2): FOXA1 is
epigenetically silenced by EZH2/PRC2 H3K27me3
deposition ($+189\%$ above normal); CDH1 is
present; the luminal programme is locked away
rather than absent. The cell is prevented from
expressing its luminal identity by an epigenetic
silencing mechanism.

The two subtypes therefore sit at opposite poles
of the luminal identity axis: ILC over-expresses
it structurally, TNBC is locked out of it
epigenetically. Both are derangements of luminal
identity --- from opposite directions.

The therapeutic consequence is direct and inverted:
ILC requires working \textit{with} the hyperactivated
FOXA1/ER circuit by fully degrading the ER receptor
(fulvestrant $>$ aromatase inhibitor), because
FOXA1 amplification makes ligand depletion alone
insufficient. TNBC requires first removing the
EZH2 silencing lock (tazemetostat) to restore FOXA1
and ESR1 before any endocrine agent can engage.

The individual features are each published in
isolation. The unifying geometric inversion framing
--- that ILC and TNBC occupy opposite poles of the
same axis with mechanistically opposite locks and
therapeutically opposite entry points --- has no
published equivalent. Verdict: CONVERGENT-NOVEL.
\end{abstract}

\vspace{0.5em}
\noindent\textbf{Keywords:}
ILC; invasive lobular carcinoma; TNBC;
triple-negative breast cancer; FOXA1; EZH2;
CDH1; E-cadherin; geometric inverse; structural lock;
epigenetic lock; luminal identity axis; fulvestrant;
aromatase inhibitor; tazemetostat; Waddington landscape;
attractor geometry; OrganismCore; GSE176078; Type~3;
H3K27me3; PRC2; CONVERGENT-NOVEL

\newpage
\tableofcontents
\newpage

% ══════════════════════════════════════════════════════════════
\section{Document Metadata}
% ══════════════════════════════════════════════════════════════

\begin{center}
\begin{tabular}{ll}
\toprule
\textbf{Field} & \textbf{Value} \\
\midrule
Document ID & CS-LIT-7 \\
Parent document & BRCA-S8h (2026-03-05) \\
Verdict & CONVERGENT-NOVEL \\
Date & 2026-03-06 \\
Author & Eric Robert Lawson \\
ORCID & 0009-0002-0414-6544 \\
Organisation & OrganismCore \\
Contact & OrganismCore@proton.me \\
Repository &
\href{https://github.com/Eric-Robert-Lawson/attractor-oncology}
{attractor-oncology} (ID: 1172048282) \\
Derivation data & GSE176078 (19,542 single cancer cells) \\
CS-LIT-1 DOI &
\href{https://doi.org/10.5281/zenodo.18883922}
{10.5281/zenodo.18883922} \\
CS-LIT-2 DOI &
\href{https://doi.org/10.5281/zenodo.18884158}
{10.5281/zenodo.18884158} \\
CS-LIT-16 DOI &
\href{https://doi.org/10.5281/zenodo.18884003}
{10.5281/zenodo.18884003} \\
CS-LIT-22 DOI &
\href{https://doi.org/10.5281/zenodo.18892788}
{10.5281/zenodo.18892788} \\
License & CC BY 4.0 \\
Biological contradictions & 0 \\
\bottomrule
\end{tabular}
\end{center}

% ════���═════════════════════════════════════════════════════════
\section{The Waddington Framework: Attractor
Lock Types}
% ══════════════════════════════════════════════════════════════

The OrganismCore framework applies Waddington landscape
geometry to cancer. In this model, each stable cancer
state corresponds to a specific attractor geometry with
a specific molecular maintenance mechanism (the lock)
and a specific therapeutic action required to dissolve
it.

Four canonical attractor types are identified across
breast cancer:

\begin{center}
\begin{tabular}{p{1.5cm}p{2.5cm}p{8cm}}
\toprule
\textbf{Type} & \textbf{Name} & \textbf{Description} \\
\midrule
Type~1 &
Slope arrest &
Cell arrested on the luminal differentiation slope.
Identity programme partially or fully active.
A molecular brake prevents completion of normal
differentiation. \\[0.4em]
Type~2 &
Wrong valley &
Cell fallen into the wrong attractor.
Correct identity programme epigenetically silenced.
A lock maintains the wrong state. \\[0.4em]
Type~3 &
Structural lock &
Cell held in an arrested state by loss of a
structural component (cell adhesion, cytoskeletal)
that normally constrains cell identity. \\[0.4em]
Type~4 &
Root lock &
Cell arrested at the progenitor root.
No terminal identity programme present ---
it was never activated, not silenced. \\
\bottomrule
\end{tabular}
\end{center}

ILC is \textbf{Type~3 (structural lock)}. TNBC is
\textbf{composite Type~1$\rightarrow$2 (epigenetic
lock)}. They are the only two subtypes that sit at
the extremes of the luminal identity axis --- and
they sit there by opposite mechanisms.

% ══════════════════════════════════════════════════════════════
\section{The Luminal Identity Axis}
% ══════════════════════════════════════════════════════════════

\subsection{What the axis measures}

The luminal identity axis is defined by the
FOXA1/EZH2 ratio: the balance between the pioneer
transcription factor that opens and maintains luminal
chromatin (FOXA1) and the Polycomb repressor that
silences it (EZH2). This axis orders all five
measurable breast cancer subtypes across seven
independent datasets comprising ${\sim}7{,}500$
patients (CS-LIT-1,
\href{https://doi.org/10.5281/zenodo.18883922}
{DOI: 10.5281/zenodo.18883922}).

\subsection{Where ILC and TNBC sit on the axis}

\begin{center}
\colorbox{axisbox}{
\begin{minipage}{0.88\textwidth}
\vspace{0.5em}
\centering
\textbf{THE LUMINAL IDENTITY AXIS}\\[0.6em]
\small
$\underbrace{\text{CL}(0.10)}_{\text{Root lock}}$
\hspace{0.5em}$\leftarrow$\hspace{0.5em}
$\underbrace{\text{TNBC}(0.52)}_{\text{Epigenetic lock}}$
\hspace{0.5em}$\leftarrow$\hspace{0.5em}
$\underbrace{\text{HER2}(3.34)}_{\text{Kinase override}}$
\hspace{0.5em}$\leftarrow$\hspace{0.5em}
$\underbrace{\text{LumB}(8.10)}_{\text{Chromatin lock}}$
\hspace{0.5em}$\leftarrow$\hspace{0.5em}
$\underbrace{\text{LumA}(9.38)}_{\text{Cycle lock}}$
\hspace{0.5em}$\leftarrow$\hspace{0.5em}
$\underbrace{\textcolor{green!50!black}{\text{ILC}}}_{\text{FOXA1 }\uparrow\uparrow\text{ above normal}}$\\[0.6em]
\normalsize
Values are FOXA1/EZH2 ratios from GSE176078
(scRNA-seq, per-patient means).\\
ILC sits \textit{above} LumA --- hyperactivated
beyond the normal luminal maximum.\\
TNBC sits at 0.52 --- FOXA1 epigenetically
silenced by EZH2.\\[0.3em]
\textbf{ILC and TNBC are the opposite poles.
The axis between them is the full range of
luminal identity expression in breast cancer.}
\vspace{0.5em}
\end{minipage}}
\end{center}

% ══════════════════════════════════════════════════════════════
\section{ILC: The Structural Lock}
% ══════════════════════════════════════════════════════════════

\subsection{The molecular architecture}

\begin{center}
\colorbox{ilcbox}{
\begin{minipage}{0.88\textwidth}
\vspace{0.5em}
\textbf{ILC ATTRACTOR STATE --- TYPE~3
STRUCTURAL LOCK}\\[0.4em]
\textbf{FOXA1:} Hyperactivated above normal
LumA levels. The luminal pioneer factor is
overdriving the ER circuit --- more ER activity
than normal, not less.\\[0.3em]
\textbf{EZH2:} Low. The epigenetic silencing
programme is not engaged. This is the opposite
of TNBC.\\[0.3em]
\textbf{CDH1 (E-cadherin):} Absent.
${\sim}65\%$ of ILC cases carry CDH1 mutations;
${\sim}35\%$ have CDH1 promoter methylation.
Combined loss rate approaches 100\% in classical
ILC.\\[0.3em]
\textbf{Lock mechanism:} CDH1/E-cadherin is the
cell adhesion molecule that organises luminal
cells into coherent epithelial structures. Its
loss removes the structural constraint on cell
identity --- the cell knows it is luminal
(FOXA1 hyperactive, ER signalling amplified)
but is physically disinhibited. The lock is
structural, not epigenetic.\\[0.3em]
\textbf{Attractor type:} Type~3 --- structural
lock. Not Type~1 (the luminal programme is not
merely arrested; it is overdriven). Not Type~2
(the cell has not fallen into the wrong valley;
it is hypercommitted to luminal identity).
\vspace{0.5em}
\end{minipage}}
\end{center}

\subsection{What FOXA1 hyperactivation means
in ILC}

In normal luminal epithelial cells, FOXA1 is the
pioneer factor that opens luminal chromatin and
enables ER binding. In LumA, FOXA1 is high and
the ER circuit is active. In ILC, FOXA1 is
\textit{higher than in LumA} --- it has been
amplified beyond the normal maximum.

The consequences:

\begin{itemize}
    \item ER circuit output is amplified beyond
    normal levels --- more ER-driven transcription
    than a normal luminal cell
    \item Ligand depletion (aromatase inhibitor)
    reduces estrogen but cannot overcome FOXA1
    amplification of the ER circuit: even reduced
    ligand drives supranormal ER output through
    an overactivated pioneer factor apparatus
    \item Full receptor degradation (fulvestrant)
    degrades the ER protein itself, bypassing the
    FOXA1 amplification entirely
    \item This is why fulvestrant is predicted
    to be superior to AIs in FOXA1-high ILC ---
    not because of ER mutation (the usual rationale
    for fulvestrant preference) but because of
    FOXA1 circuit amplification
\end{itemize}

\subsection{CDH1 loss: the structural component}

CDH1 encodes E-cadherin, the primary cell adhesion
molecule in luminal breast epithelium. In normal
luminal cells, E-cadherin:

\begin{itemize}
    \item Maintains cell-cell contacts and
    epithelial sheet organisation
    \item Constrains cell morphology, preventing
    the discohesive single-file growth pattern
    characteristic of ILC
    \item Participates in contact inhibition
    signalling that limits proliferation
    \item Acts as a structural node that links
    the cytoskeleton to the extracellular
    environment
\end{itemize}

CDH1 loss in ILC removes this structural
constraint. The cell retains its luminal identity
programme (FOXA1 hyperactive) but is released from
the physical organisation that normal luminal cells
require. The result is the characteristic ILC
phenotype: discohesive, single-file infiltration
with intact luminal marker expression. The lock is
structural, not epigenetic.

% ══════════════════════════════════════════════════════════════
\section{TNBC: The Epigenetic Lock}
% ══════════════════════════════════════════════════════════════

\begin{center}
\colorbox{tnbcbox}{
\begin{minipage}{0.88\textwidth}
\vspace{0.5em}
\textbf{TNBC ATTRACTOR STATE --- COMPOSITE
TYPE~1$\rightarrow$2 EPIGENETIC LOCK}\\[0.4em]
\textbf{FOXA1:} Epigenetically silenced by
EZH2/PRC2. Near-zero expression. The luminal
pioneer factor is not absent from the genome ---
it is present but locked away by H3K27me3
deposition on the FOXA1 promoter.\\[0.3em]
\textbf{EZH2:} $+189\%$ above normal Mature
Luminal reference (highest of all subtypes).
The Polycomb repressor is maximally active,
silencing FOXA1, GATA3, and ESR1.\\[0.3em]
\textbf{CDH1:} Present. The structural
constraint is intact; it is the epigenetic
identity programme that is locked, not the
structural architecture.\\[0.3em]
\textbf{Lock mechanism:} PRC2/EZH2-mediated
H3K27me3 deposition on FOXA1, GATA3, and
ESR1 promoters silences the luminal identity
programme. The cell has the structural capacity
for luminal organisation (CDH1 present) but
is prevented from expressing luminal identity
by an epigenetic lock.\\[0.3em]
\textbf{Attractor type:} Composite
Type~1$\rightarrow$2. The cell began
on the luminal slope (Type~1) but EZH2
overexpression has driven it into an
alternative attractor state where luminal
identity is actively maintained in a silenced
configuration (Type~2 component).
\vspace{0.5em}
\end{minipage}}
\end{center}

% ══════════════════════════════════════════════════════════════
\section{The Geometric Inversion: Side by Side}
% ══════════════════════════════════════════════════════════════

\subsection{Full comparison table}

\begin{center}
\begin{tabular}{p{3.5cm}p{4cm}p{4cm}}
\toprule
\textbf{Feature} &
\textbf{\textcolor{green!50!black}{ILC}} &
\textbf{\textcolor{red!70!black}{TNBC}} \\
\midrule
FOXA1 expression &
Hyperactivated above
normal LumA &
Silenced to near-zero
by EZH2/PRC2 \\[0.4em]
EZH2 expression &
Low &
$+189\%$ above normal \\[0.4em]
CDH1 (E-cadherin) &
Absent
(${\sim}65\%$ mutation,
${\sim}35\%$ methylation) &
Present \\[0.4em]
Luminal identity &
Over-expressed;
hypercommitted &
Epigenetically silenced;
locked away \\[0.4em]
Lock type &
Structural (CDH1 loss) &
Epigenetic
(EZH2/PRC2 H3K27me3) \\[0.4em]
Attractor geometry &
Type~3 &
Composite 1$\rightarrow$2 \\[0.4em]
Direction of derangement &
Too much luminal identity,
structurally uncontained &
Luminal identity present
but epigenetically locked \\[0.4em]
ER status (clinical) &
ER-positive &
ER-negative \\[0.4em]
Mechanism of ER status &
ER hyperactive
(FOXA1 amplification) &
ER silenced
(H3K27me3 on ESR1) \\[0.4em]
Therapeutic first step &
Full ER degradation
(fulvestrant) &
EZH2 inhibition
(tazemetostat) \\[0.4em]
Why that step &
FOXA1 amplification bypasses
ligand depletion; ER protein
must be degraded directly &
ESR1 is locked by EZH2;
must remove lock before
any endocrine agent can engage \\
\bottomrule
\end{tabular}
\end{center}

\subsection{The inversion is precise, not
metaphorical}

The FOXA1/EZH2 axis is a physical measurement
of chromatin state. ILC sits at the highest
end of this axis (FOXA1 hyperactivated, EZH2
low). TNBC sits near the lowest end (FOXA1
silenced, EZH2 maximal). Both involve luminal
identity. Both involve the same proteins. The
direction is exactly inverted.

CDH1 is also inverted: absent in ILC, present
in TNBC. This is not coincidental --- CDH1 loss
is the structural mechanism that produces the
ILC phenotype. CDH1 presence in TNBC is
consistent with an epigenetically arrested
cell that retains structural organisation.

The inversion extends to therapeutic logic:
the correct first step in ILC degrades the
overactive ER receptor (fulvestrant). The
correct first step in TNBC removes the
epigenetic lock that prevents the ER from
being expressed at all (tazemetostat). Giving
fulvestrant to a TNBC patient is giving a
receptor-degrader to a patient with no receptor
to degrade. Giving an AI to an ILC patient
is depleting ligand from a system whose FOXA1
amplification makes ligand abundance irrelevant.

Both errors arise from not recognising the
geometric inversion.

% ══════════════════════════════════════════════════════════════
\section{Published Literature: Confirmation
of Individual Features}
% ══════════════════════════════════════════════════════════════

\subsection{What IS in the literature}

\begin{enumerate}
    \item \textbf{CDH1 loss in ILC:}
    Well-established. ${\sim}65\%$ CDH1 somatic
    mutation, ${\sim}35\%$ CDH1 methylation in
    classical ILC. Published across multiple
    large ILC genomic studies (TCGA ILC, Ciriello
    et al.\ \textit{Cell} 2015).

    \item \textbf{FOXA1 high in ILC:}
    Established. Frontiers Oncology 2025 ILC
    review explicitly notes FOXA1 hyperactivation
    in ILC and associates it with altered endocrine
    sensitivity. FOXA1 mutations in ILC (amplifying
    and gain-of-function) published in Cancer Cell
    2020.

    \item \textbf{FOXA1 absent in TNBC:}
    Published. FOXA1 is near-zero in TNBC/basal
    breast cancer across multiple datasets.

    \item \textbf{EZH2 high in TNBC:}
    Confirmed. Schade et al.\ \textit{Nature}
    2024 (Harvard/Ludwig): EZH2/PRC2 is the
    convergence node in basal-like TNBC that
    silences FOXA1, GATA3, and ESR1.
    Frontiers Oncology 2020: EZH2/NSD2 axis,
    ${\sim}4{,}000$ patients, EZH2 highest
    in TNBC.

    \item \textbf{CDH1 present in TNBC:}
    Consistent with published TNBC biology.
    TNBC is typically classified as basal-like
    with intact epithelial markers except for
    the loss of ER, PR, and HER2.

    \item \textbf{Fulvestrant vs AI in ILC:}
    Active area of investigation. The Frontiers
    Oncology 2025 ILC review explicitly notes
    the open question of whether FOXA1-driven
    altered endocrine sensitivity changes the
    AI vs fulvestrant choice. NCT02206984
    compares endocrine agents in ILC (Ki67
    endpoint; ongoing).
\end{enumerate}

\subsection{What is NOT in the literature}

\begin{center}
\colorbox{novelbox}{
\begin{minipage}{0.88\textwidth}
\vspace{0.5em}
\textbf{The following have no published equivalent
as of 2026-03-06:}
\begin{itemize}[noitemsep]
    \item The framing of ILC and TNBC as
    \textit{geometric inverses} on the same
    luminal identity axis.
    \item The statement: ``ILC = FOXA1
    hyperactivated above normal, CDH1 absent,
    structural lock; TNBC = FOXA1 epigenetically
    silenced by EZH2, CDH1 present, epigenetic
    lock; these are opposite poles of the same
    axis.''
    \item The derivation that the therapeutic
    entry points are mechanistically inverted:
    fulvestrant as the \textit{correct} first
    step in ILC (FOXA1 amplification bypasses
    ligand depletion); tazemetostat as the
    \textit{correct} first step in TNBC
    (epigenetic lock must be dissolved before
    endocrine therapy can engage).
    \item The explicit mechanistic explanation
    for fulvestrant superiority over AI in ILC
    grounded in FOXA1 circuit amplification
    (not ESR1 mutation, which is the conventional
    rationale).
    \item No published paper uses this geometric
    inversion frame. No published paper derives
    treatment logic from it.
\end{itemize}
\vspace{0.5em}
\end{minipage}}
\end{center}

% ══════════════════════════════════════════════════════════════
\section{Therapeutic Consequences}
% ══════════════════════════════════════════════════════════════

\subsection{ILC: why fulvestrant, not AI}

The standard endocrine therapy debate in ILC is
framed around FOXA1 mutations (gain-of-function
mutations drive ligand-independent ER activation,
making AIs theoretically less effective). The
geometric framework provides an additional, more
general reason that applies to all FOXA1-high ILC
regardless of mutation status:

FOXA1 is the pioneer factor that opens chromatin
for ER binding. When FOXA1 is hyperactivated
(elevated above normal LumA), the ER circuit is
primed to maximal sensitivity --- it operates at
gain $>1$. In this state:

\begin{itemize}
    \item AI reduces estrogen ligand. Ligand
    reduction suppresses ER signalling proportionally
    in a normal luminal cell. In a FOXA1-amplified
    ILC cell, the ER circuit amplification partially
    compensates for reduced ligand --- the gain
    in the FOXA1 component maintains circuit
    activity despite lower input signal.
    \item Fulvestrant degrades the ER protein
    itself. This acts upstream of FOXA1: if there
    is no ER protein, FOXA1 amplification is
    irrelevant because there is no receptor to
    amplify. The circuit is interrupted at the
    receptor level.
\end{itemize}

The prediction: in FOXA1-high ILC, fulvestrant
produces deeper suppression of ER circuit output
than AI because it eliminates the receptor that
FOXA1 amplification would otherwise maintain
active. The greater the FOXA1 hyperactivation,
the greater the advantage of fulvestrant over AI.

This prediction is the subject of CS-LIT-18
(companion document) and is testable in
NCT02206984 tissue archives.

\subsection{TNBC: why tazemetostat before
fulvestrant}

The tazemetostat $\rightarrow$ fulvestrant
sequence for TNBC is fully documented in
CS-LIT-16
(\href{https://doi.org/10.5281/zenodo.18884003}
{DOI: 10.5281/zenodo.18884003}). The relevant
comparison here is with ILC:

In TNBC, giving fulvestrant first makes no
sense because there is no expressed ER protein
to degrade. The ESR1 gene is present in the
genome but its promoter is locked by H3K27me3.
The sequence must be tazemetostat first (removes
H3K27me3 from FOXA1 and ESR1 promoters),
then FOXA1 re-expresses, then ESR1 is
transcribed, then ER protein is produced, and
only then can fulvestrant engage.

In ILC, the opposite is true: the ER is
hyperactive. Fulvestrant is the correct first
step precisely because there is too much ER
activity, not too little.

The same drug (fulvestrant) is the correct
last step in the TNBC sequence and the correct
first step in the ILC approach --- but for
completely inverted biological reasons.

% ══════════════════════════════════════════════════════════════
\section{Position in the Published Series}
% ══════════════════════════════════════════════════════════════

CS-LIT-7 is the foundational architecture
document for ILC within the OrganismCore BRCA
cross-subtype series. It must precede CS-LIT-18
(fulvestrant superiority in ILC, companion
document, currently in preparation).

\begin{center}
\begin{tabular}{p{2.8cm}p{8.5cm}}
\toprule
\textbf{Document} & \textbf{Role} \\
\midrule
CS-LIT-7\\(this document) &
The geometric architecture: ILC as the structural
inverse of TNBC on the luminal identity axis.
The mechanistic foundation for all ILC treatment
logic in this framework. \\[0.5em]
CS-LIT-18\\(in preparation) &
The clinical prediction: fulvestrant superiority
over AIs in FOXA1-stratified ILC. Depends on
CS-LIT-7 as its mechanistic foundation. \\[0.5em]
CS-LIT-1\\
\href{https://doi.org/10.5281/zenodo.18883922}
{DOI: 10.5281/zenodo.18883922} &
The FOXA1/EZH2 ratio ordering axis. ILC sits
above LumA on this axis (FOXA1 hyperactivated);
TNBC at 0.52. CS-LIT-7 explains why. \\[0.5em]
CS-LIT-2\\
\href{https://doi.org/10.5281/zenodo.18884158}
{DOI: 10.5281/zenodo.18884158} &
The six lock type classification. ILC is Lock
Type~3 (structural). TNBC is composite
Type~1$\rightarrow$2. CS-LIT-7 is the
dedicated evidence document for this
classification of both subtypes. \\[0.5em]
CS-LIT-16\\
\href{https://doi.org/10.5281/zenodo.18884003}
{DOI: 10.5281/zenodo.18884003} &
The tazemetostat $\rightarrow$ fulvestrant
sequence for TNBC. The therapeutic mirror
image of CS-LIT-18. \\[0.5em]
CS-LIT-22\\
\href{https://doi.org/10.5281/zenodo.18892788}
{DOI: 10.5281/zenodo.18892788} &
FOXA1/EZH2 dual IHC as point-of-care tool.
The IHC instrument that places ILC above LumA
and TNBC below LumA on the clinical slide. \\
\bottomrule
\end{tabular}
\end{center}

% ═══════════════════════════════════════════════════════════���══
\section{Locked Statement}
% ══════════════════════════════════════════════════════════════

\begin{center}
\colorbox{findingbox}{
\begin{minipage}{0.88\textwidth}
\vspace{0.6em}
\textbf{Stated once, precisely, without
inflation:}\\[0.5em]
ILC and TNBC occupy opposite poles of the luminal
identity axis. ILC is a Type~3 structural lock:
FOXA1 hyperactivated above normal LumA levels,
EZH2 low, CDH1 absent (${\sim}65\%$ mutation,
${\sim}35\%$ methylation). TNBC is a composite
Type~1$\rightarrow$2 epigenetic lock: FOXA1
epigenetically silenced by EZH2/PRC2
($+189\%$ above normal), CDH1 present. Both are
derangements of luminal identity from opposite
directions: ILC over-expresses it structurally,
TNBC is locked out of it epigenetically. The
therapeutic entry points are therefore inverted:
ILC requires full ER receptor degradation
(fulvestrant $>$ AI) because FOXA1 amplification
bypasses ligand depletion; TNBC requires EZH2
inhibition (tazemetostat) to restore FOXA1 and
ESR1 before any endocrine agent can engage. The
individual features (FOXA1-high ILC, FOXA1-absent
TNBC, CDH1-absent ILC, CDH1-present TNBC, EZH2-high
TNBC, EZH2-low ILC) are each published in isolation.
The unifying geometric inversion framing --- same
axis, opposite poles, opposite locks, opposite
therapeutic entry points --- has no published
equivalent as of 2026-03-06. Verdict:
CONVERGENT-NOVEL.
\vspace{0.5em}
\end{minipage}}
\end{center}

\end{document}
